% ============================================================
% Come si usano, argomento 4: condizionamento e indipendenza
% ============================================================
\argomento{Probabilità condizionata e indipendenza}{Formule a p.~\pageref*{f:condizionata}}[Condizionamento e indipendenza]

\begin{scelta}{Come scegliere la regola: dal testo alla formula, e l'albero}
\renewcommand{\arraystretch}{1.25}
\begin{tabular}{@{}P{0.5\linewidth}P{0.44\linewidth}@{}}
  \toprule
  \textbf{Se il testo\ldots} & \textbf{si usa\ldots} \\
  \midrule
  dice «sapendo che\ldots», «dato che\ldots», «se si è verificato\ldots»
    & la probabilità condizionata \vv{condizionata} \\
  chiede che succedano due cose, in due tempi (estrazioni successive)
    & la probabilità composta \vv{composta}: si moltiplica lungo un ramo \\
  dà la probabilità dell'evento in alcuni «casi» (quale bit è stato trasmesso, quale urna, malato o sano)
    & la probabilità totale \vv{totale}: si sommano i rami \\
  osserva l'effetto e chiede la causa («ricevuto 1: era stato trasmesso 1?»)
    & Bayes \vv{bayes}: un ramo diviso la somma dei rami \\
  dice che le prove non si influenzano, «indipendenti»
    & il prodotto delle probabilità \vv{indipendenza} \\
  \bottomrule
\end{tabular}

Le tre leggi si leggono tutte sullo stesso albero. Esempio: canale binario, $E_1$ = «trasmesso 0», $E_2$ =
«trasmesso 1», $A$ = «ricevuto 1».
\begin{center}
\begin{tikzpicture}[x=1cm, y=0.52cm, every node/.style={font=\small},
                    nodo/.style={draw, rounded corners=2pt, inner sep=2.5pt, fill=white}]
  \node[draw, circle, fill=mspilight, inner sep=2pt] (r) at (0,0) {};
  \node[nodo, fill=mspiA!12] (t0) at (2.7,1.8) {$E_1$: 0 tr.};
  \node[nodo, fill=mspiA!12] (t1) at (2.7,-1.8) {$E_2$: 1 tr.};
  \draw (r) -- node[above, sloped, font=\scriptsize] {$0{,}6$} (t0);
  \draw (r) -- node[below, sloped, font=\scriptsize] {$0{,}4$} (t1);
  \node[nodo] (a) at (6.0,2.8) {0 ric.};
  \node[nodo, fill=evid] (b) at (6.0,0.9) {$A$: 1 ric.};
  \node[nodo] (c) at (6.0,-0.9) {0 ric.};
  \node[nodo, fill=evid] (d) at (6.0,-2.8) {$A$: 1 ric.};
  \draw (t0.east) -- node[above, sloped, font=\scriptsize] {$0{,}9$} (a.west);
  \draw (t0.east) -- node[below, sloped, font=\scriptsize] {$0{,}1$} (b.west);
  \draw (t1.east) -- node[above, sloped, font=\scriptsize] {$0{,}1$} (c.west);
  \draw (t1.east) -- node[below, sloped, font=\scriptsize] {$0{,}9$} (d.west);
  \node[anchor=west] at (7.0,0.9) {$0{,}6 \cdot 0{,}1 = 0{,}06$};
  \node[anchor=west] at (7.0,-2.8) {$0{,}4 \cdot 0{,}9 = 0{,}36$};
  \node[anchor=west, text width=7.4cm, align=flush left, font=\small] at (9.9,0)
    {\textbf{composta} \vv{composta}: un percorso è il prodotto dei suoi rami,
     $P(A \cap E_2) = 0{,}4 \cdot 0{,}9 = 0{,}36$.\\[4pt]
     \textbf{totale} \vv{totale}: si sommano i percorsi che finiscono in $A$,
     $P(A) = 0{,}06 + 0{,}36 = 0{,}42$.\\[4pt]
     \textbf{Bayes} \vv{bayes}: un percorso diviso la somma,
     ${P(E_2 \mid A)} = 0{,}36 / 0{,}42 = 6/7$.};
\end{tikzpicture}
\end{center}
\end{scelta}

\scheda{condizionata}{
  \usoquando{Il testo dice «sapendo che», «dato che», «se si è verificato»: c'è un'informazione parziale sul
    risultato.}
  \usocome{Si restringe il mondo all'evento noto $B$ e si guarda quale parte di $B$ è occupata da $A$. In uno
    spazio equiprobabile basta contare dentro $B$.}
  \usoattenzione{$P(A \mid B)$ e $P(B \mid A)$ sono numeri diversi: l'evento noto va al denominatore. Serve
    $P(B) > 0$.}
  \usoesempio{Un dado onesto: $P(6 \mid \text{pari}) = \frac{1/6}{3/6} = \frac{1}{3}$;
    $P(\text{dispari} \mid \{1, 2, 3\}) = \frac{2/6}{3/6} = \frac{2}{3}$; invece $P(\text{pari} \mid 6) = 1$.}
}

\scheda{composta}{
  \usoquando{Si chiede che succedano due o più cose, e l'esperimento procede per tempi successivi in cui ciò che
    accade dopo dipende da ciò che è accaduto prima: estrazioni senza reinserimento, prove a cascata.}
  \usocome{La probabilità della prima cosa per la probabilità della seconda \emph{sapendo} che la prima è successa.
    Sull'albero si moltiplicano le probabilità lungo il percorso.}
  \usovarianti{Con tre eventi la catena continua: ogni fattore è condizionato a tutti i precedenti.}
  \usoesempio{Due assi di fila da 52 carte: $\frac{4}{52} \cdot \frac{3}{51} = \frac{1}{221} \simeq 0{,}0045$. Tre
    assi di fila: $\frac{4}{52} \cdot \frac{3}{51} \cdot \frac{2}{50} = \frac{1}{5\,525}$.}
}

\scheda{partizione}{
  \usoquando{Per impostare la probabilità totale e il teorema di Bayes: sono i «casi» in cui si divide il mondo
    (quale bit è stato trasmesso, malato o sano, quale urna è stata scelta).}
  \usoattenzione{I casi devono coprire tutto e non sovrapporsi. Controllo rapido: le loro probabilità sommano a 1.}
}

\scheda{totale}{
  \usoquando{La probabilità di $A$ non si sa calcolare direttamente, ma si sa calcolare in ciascuno dei casi
    $E_i$.}
  \usocome{(1) Si sceglie la partizione. (2) Si scrivono le $P(E_i)$ e le $P(A \mid E_i)$. (3) Si somma
    $P(A \mid E_i)\,P(E_i)$ su tutti i casi: sull'albero, la somma dei percorsi che finiscono in $A$.}
  \usoesempio{Canale binario: $P(1\ \text{ric.}) = 0{,}1 \cdot 0{,}6 + 0{,}9 \cdot 0{,}4 = 0{,}42$.}
}

\scheda{bayes}{
  \usoquando{Si osserva l'effetto e si chiede la causa: si conosce $P(\text{effetto} \mid \text{causa})$ e si vuole
    $P(\text{causa} \mid \text{effetto})$.}
  \usocome{Al numeratore il percorso che passa per la causa che interessa; al denominatore la probabilità totale
    dell'effetto, cioè la somma di tutti i percorsi che ci arrivano.}
  \usoattenzione{Conta anche la probabilità a priori $P(E_i)$: un test affidabile per una malattia rara dà molti
    falsi positivi.}
  \usoesempio{Canale: $P(1\ \text{tr.} \mid 1\ \text{ric.}) = \frac{0{,}9 \cdot 0{,}4}{0{,}42} = \frac{6}{7} \simeq
    0{,}857$. Test diagnostico (malattia nell'1\% della popolazione, test positivo nel 99\% dei malati e nel 5\% dei
    sani): $P(M \mid +) = \frac{0{,}99 \cdot 0{,}01}{0{,}99 \cdot 0{,}01 + 0{,}05 \cdot 0{,}99} =
    \frac{0{,}0099}{0{,}0594} \simeq 0{,}167$.}
}

\scheda{indipendenza}{
  \usoquando{Per \emph{calcolare}: se il testo dice che le prove sono indipendenti (lanci, bit, estrazioni con
    reinserimento), la probabilità che succedano tutte le cose è il prodotto. Per \emph{verificare}: si controlla se
    $P(A \cap B) = P(A)\,P(B)$.}
  \usovarianti{Se $A$ e $B$ sono indipendenti lo sono anche i complementari: utile con «nessuno», insieme a De
    Morgan (\vv{demorgan}).}
  \usoattenzione{Indipendenti non vuol dire incompatibili: due eventi incompatibili con probabilità positive sono
    \emph{dipendenti}, perché $P(A \cap B) = 0 \neq P(A)\,P(B)$.}
  \usoesempio{Un lancio di dado, $A$ = «pari», $B = \{1, 2\}$: $P(A \cap B) = \frac{1}{6} = \frac{1}{2} \cdot
    \frac{1}{3}$, quindi sono indipendenti. Due bit indipendenti, ciascuno sbagliato con probabilità $0{,}1$:
    «nessuno sbagliato» ha probabilità $0{,}9 \cdot 0{,}9 = 0{,}81$.}
}
